\section{Method}
\label{sec:methodology}

\subsection{Material Conversion as the Intervention}
\label{sec:conversion_intervention}

The intervention stays within the same USD scene family. ConvertAsset starts
from a root USD file, recursively discovers referenced layers, payloads,
variants, and clips, and writes \texttt{\_noMDL} sibling USD assets; it does not
flatten the scene. Converted arcs target those siblings; sources stay read-only.

Within each processed USD file, MDL-driven materials are replaced by a
core-channel PreviewSurface approximation rather than by a full MDL semantic
compiler. ConvertAsset builds portable base-color, roughness, metallic, and
normal networks; texture inputs are resolved relative to their declaring layer
when possible, otherwise constant fallbacks are used. This deliberately leaves
many MDL mechanisms outside the conversion target. Clearcoat, procedural
textures, displacement or height cues, opacity, emission, and other specialized
mechanisms therefore become risk labels for downstream audit, not promises that
the mechanism has been faithfully reproduced.

\subsection{Evidence Gates}
\label{sec:evidence_gates}

We organize the evaluation as four evidence gates. The first gate measures
paired proxy similarity on four Isaac Sim furniture assets: PSNR, SSIM, LPIPS,
CLIP similarity, DINOv2 similarity, and headless load/render behavior. This
gate checks that the conversion can preserve obvious visual structure in a
controlled single-object setting.

The second gate is the VLM grounding protocol on GRScenes. We use
the original GRScenes tree as the source condition, materialize noMDL
counterparts in a scratch workspace, render matched target-centered views, and
project target bounding boxes into each image. Visual QA separates clean
PASS-only pairs from target-visible PASS/WARN material-shift pairs. Local VLM
backends are prompted to return a target-category answer and a point location:
an Unsloth-quantized Gemma4 checkpoint
\texttt{unsloth/\allowbreak gemma-4-E4B-it-\allowbreak unsloth-bnb-4bit} is
the canonical run, while
\texttt{Qwen/\allowbreak Qwen2.5-VL-7B-\allowbreak Instruct} is a second-model
diagnostic under the same manifest. We score answer match, point-in-box under
normalized-1000 coordinates, and a raw pixel-space diagnostic because coordinate
semantics proved model-dependent; these backends are measurement probes, not a
model leaderboard.

The third gate audits material mechanisms and the NVIDIA baseline. For the
selected GRScenes material-effect pool, we compare original MDL, ConvertAsset
noMDL, and NVIDIA Asset Converter PreviewSurface outputs when all three
conditions pass static gates. These static gates check condition availability,
zero active MDL in converted outputs, and target renderability; they do not by
themselves certify material-effect fidelity. For missing material bins, we use
bounded supplemental wrapper fixtures, but only as selected failure or
limitation evidence.

The fourth gate checks embodied-data usability. We run the official
InteriorAgent / KuJiaLe \texttt{val\_unseen} route through
InternNav/DualVLN~\cite{Wei2026Ground} on matched original and noMDL scenes,
and separately measure repeated official-scene load/render stability. These
runs are treated as scoped sanity evidence: they show that the converted scenes
can be exercised by a real embodied stack, but they are not broad navigation
benchmark evidence.
Table~\ref{tab:acl_evidence_gate_registry} summarizes the allowed claim
boundary for each gate.

\begin{table*}[t]
\centering
\footnotesize
\setlength{\tabcolsep}{3pt}
\begin{tabular}{p{0.15\textwidth}p{0.30\textwidth}p{0.24\textwidth}p{0.22\textwidth}}
\toprule
Gate & Evidence used & Claim supported & Forbidden promotion \\
\midrule
Proxy similarity & Four Isaac Sim assets; 24 matched render pairs; PSNR, SSIM, LPIPS, CLIP, DINOv2, and load/render diagnostics. & Converted assets pass a visual and feature-screening gate. & Task preservation, broad robustness, or speedup. \\
VLM grounding & GRScenes 15-pair clean pilot and frozen 30-pair target-centered stress set with answer, point-in-box, and coordinate scoring. & Material perturbation can be evaluated as bounded answer and point-grounding evidence. & Population-level GRScenes benchmark claims or general VLM rankings. \\
Material mechanisms & Selected opacity/transparency, emission, normal/bump, and displacement/height bins, plus clearcoat and procedural supplemental cases; original/noMDL/NVIDIA conditions are separated. & Covered bins have selected static and qualitative support; clearcoat and procedural texture remain scoped failure or limitation cases. & NVIDIA population failure rates, official-scene NVIDIA performance, or procedural-texture success. \\
Embodied-data sanity & Official KuJiaLe InternNav 99 paired episodes over three scenes plus official-scene load/render stability. & Converted scenes can enter the selected official downstream stack and load/render repeatedly. & Broad embodied-navigation robustness, all-scene preservation, or noMDL speedup. \\
\bottomrule
\end{tabular}
\caption{Evidence-gate registry for the ACL/ARR claim boundary. Each gate lists the evidence that may be used, the claim it supports, and the promotion that remains forbidden without new evidence.}
\label{tab:acl_evidence_gate_registry}
\end{table*}


\subsection{Claim Registry}
\label{sec:claim_registry}

Every major result is tied to a manifest, table, or claim entry in the project
evidence registry. The registry records not only the supported claim, but also
its scope and forbidden promotion. This is part of the method: we use conversion
as a controlled perturbation, then prevent proxy metrics, selected videos, or
single-scene results from being promoted into broader robustness claims.
